\section{Perturbation-based explanations}
\label{sec:explain}

Model-agnostic attribution needs only a real-valued function of the input, and with a class
distribution it is immediate: $f(z) = \mu_{x(z)}(\{c^\star\})$, Jev's probability of the class it
predicts on the full input when only the units in the coalition $z$ are kept (a feature line for
tabular rows, a word for text). We ran KernelSHAP~\citep{lundberg2017shap} and
LIME~\citep{ribeiro2016lime} on \ExExamplesGerman{} seeded examples per dataset, on live Jev, and
tested faithfulness by deleting the top-attributed units and re-querying.

\begin{table}[t]
  \caption{Explanations of Jev's predictions: mean drop in the predicted-class probability and
    share of label flips after deleting the top three attributed units, against three random
    units; Kendall's $\tau$ and top-3 overlap between KernelSHAP and LIME; and Jev requests used.}
  \label{tab:explain}
  \centering\small
  \setlength{\tabcolsep}{4pt}
  \fittable{\begin{tabular}{llrcccccccr}
\toprule
 & & & \multicolumn{3}{c}{Probability drop, $k=3$} & \multicolumn{2}{c}{Label flips, $k=3$} & \multicolumn{2}{c}{Agreement} & \\
\cmidrule(lr){4-6}\cmidrule(lr){7-8}\cmidrule(lr){9-10}
Dataset & Unit & Units & SHAP & LIME & random & SHAP & random & $\tau$ & top-3 & Jev requests \\
\midrule
German Credit & feature & 20 & 0.577 & 0.603 & 0.008 & 0.85 & 0.06 & 0.71 & 0.80 & 7{,}164 \\
Breast Cancer Wisconsin & feature & 30 & 0.065 & 0.155 & 0.010 & 0.05 & 0.02 & 0.37 & 0.35 & 7{,}969 \\
SST-2 & word & 18 & 0.401 & 0.419 & 0.096 & 0.40 & 0.09 & 0.51 & 0.70 & 6{,}666 \\
AG News & word & 34 & 0.075 & 0.091 & 0.003 & 0.10 & 0.00 & 0.19 & 0.27 & 7{,}919 \\
\bottomrule
\end{tabular}
}
\end{table}

Both methods run unmodified, and how useful they are depends on the dataset
(Table~\ref{tab:explain}). On German Credit
and SST-2 the explanations are faithful: deleting the three top KernelSHAP units lowers the
predicted-class probability by \ExDropShapGermanKThree{} and \ExDropShapSstKThree, against
\ExDropRandomGermanKThree{} and \ExDropRandomSstKThree{} for three random units. On Breast Cancer
Wisconsin and AG News they are weak (\ExDropShapWdbcKThree{} and \ExDropShapAgnewsKThree): with
\ExMeanUnitsWdbc{} correlated measurements, or \ExMeanUnitsAgnews{} words, no three units carry
the decision. The base value is informative in itself: with every feature removed, Jev still
gives \emph{bad} a probability of \ExEmptyInputPositiveGerman{} on German Credit, the bias that
Section~\ref{sec:metrics} measures. These explanations cost \ExRequestsTotal{} requests;
Appendix~\ref{app:explain} shows the deletion curves (Figure~\ref{fig:faithfulness}) and
example attributions (Figure~\ref{fig:attributions}).
